\documentclass{beamer}

\input{settings.tex}


\title{Luenberger Observer}
\subtitle{Control Theory, Lecture 8}
\author{by Sergei Savin}
\centering
\date{\mydate}



\begin{document}
\maketitle



%\begin{frame}{Content}
%\begin{itemize}
%\item Measurement
%\item State Estimation
%\item Observer
%\item Observation and Control
%\item Separation principle
%\end{itemize}
%\end{frame}




\begin{frame}{Measurement and control}
\begin{flushleft}

Before we considered systems and control laws of the following type:

\begin{equation}
\begin{cases}
\dot {\bo{x}} = \bo{A} \bo{x} + \bo{B} \bo{u}\\
\bo{u} = -\bo{K} \bo{x}
\end{cases}
\end{equation}

But when implementing this control law, how do we know the current value of $\bo{x}$?

\bigskip

In practice, we can \emph{estimate} it using \emph{measurement}.

\end{flushleft}
\end{frame}

\begin{frame}{Why information is imperfect?}
\begin{flushleft}

There are a number of reasons why we can not directly measure the state of the system. Here are some:

\begin{itemize}
\item \textbf{Lack of sensors.}
\item Digital measurements are discrete (available only at discrete sampling times).
\item Measurements may be noisy, delayed, or temporarily unavailable.
\item Model inaccuracies (flexibility, backlash, friction, etc.) may cause the true physical state to differ from the model state.
\item Imprecise, nonlinear and biased sensors.
\item Other physical effects.
\end{itemize}

\end{flushleft}
\end{frame}

\begin{frame}{Measurement and estimation}
\begin{flushleft}

Let us introduce an \emph{observer}. A linear system with an observer takes the following form:

\begin{equation}
\begin{cases}
\dot {\bo{x}} = \bo{A} \bo{x} + \bo{B} \bo{u} \\
\bo{y} = \bo{C} \bo{x} \\
\dot{\hat{\bo{x}}} = \text{Observer} \ (\hat{\bo{x}}, \bo{u}, \bo{y}) \\
\bo{u} = -\bo{K}\hat{\bo{x}}
\end{cases}
\end{equation}

where:

\begin{itemize}
\item $\bo{x}$ and $\bo{y}$ are the state and output (actual, true)
\item $\hat{\bo{x}}$ and $\hat{\bo{y}} =\bo{C} \hat{\bo{x}}$ are the estimated (observed) state
and output.
\end{itemize}

Notice that the controller does not know the true state $\bo{x}$, and therefore, for the control purposes, we have to use the estimated state $\hat{\bo{x}}$.

\end{flushleft}
\end{frame}




\begin{frame}{Estimation error}
	\begin{flushleft}
		
		We define state estimation error:
		
		\begin{equation}
		\bo{e} = \bo{x} - \hat{\bo{x}}
		\end{equation}
		
		Note that a controller cannot compute it, since it does not know $\bo{x}$. However, it can compare the measured output $\bo{y}$ with the estimated output $\hat{\bo{y}} =\bo{C} \hat{\bo{x}}$:
		
		\begin{equation}
		\Tilde{\bo{y}} = \bo{y} -  \bo{C} \hat{\bo{x}}  
		\end{equation}		
		
		This can always be computed.
		
	\end{flushleft}
\end{frame}



\begin{frame}{Estimation - dynamics}
\begin{flushleft}

Let us consider autonomous dynamical system
\begin{equation}
\label{eq:LTI}
\begin{cases}
\dot {\bo{x}} = \bo{A} \bo{x} + \bo{B} \bo{u} \\
\bo{y} = \bo{C} \bo{x}
\end{cases}
\end{equation}
%
with measurements $\bo{y}$. We want to get as good an estimate of the state $\hat{\bo{x}}$ as we can.

\bigskip

Proposal: the dynamics should also hold for the observed state:
\begin{equation}
\dot{\hat{\bo{x}}} = \bo{A} \hat{\bo{x}} + \bo{B} \bo{u}
\end{equation}
%
If we knew the initial conditions exactly, we could find the exact state of the system without using measurements $\bo{y}$. We can call it an open loop observation. Unfortunately, we do not know the initial conditions precisely.


\end{flushleft}
\end{frame}





\begin{frame}{Estimation - observer}
\begin{flushleft}

We propose \emph{observer} that takes into account measurements in a linear way; similar to linear control $-\bo{K}\bo{x}$, here we propose a linear correction law $\bo{L}\Tilde{\bo{y}}$. Remembering that $\Tilde{\bo{y}} = \bo{y} -  \bo{C} \hat{\bo{x}} $ we get:

\begin{equation}
\label{eq:Observer}
\dot{\hat{\bo{x}}}  = \bo{A} \hat{\bo{x}} + \bo{B} \bo{u} + \bo{L}(\bo{y} - \bo{C} \hat{\bo{x}})
\end{equation}
%
This is called \emph{Luenberger observer}. The next task is to find suitable observer gain $ \bo{L}$.

\bigskip


We subtract from $\dot {\bo{x}} = \bo{A} \bo{x} + \bo{B} \bo{u}$ the observer \eqref{eq:Observer}, to get the \emph{observer error dynamics}:

\begin{align}
	\dot {\bo{x}} - \dot{\hat{\bo{x}}} = 
	\bo{A} \bo{x} - \bo{A} \hat{\bo{x}} - 
	\bo{L}(\mathbf y - \bo{C} \hat{\bo{x}})
	\\
	\dot {\bo{x}} - \dot{\hat{\bo{x}}}
	= 
	\bo{A} (\bo{x} -   \hat{\bo{x}} ) - 
	\bo{L}(\bo{C} \bo{x} -  \bo{C}\hat{\bo{x}}  ) 
\end{align}
%
\begin{equation}
\dot {\bo{e} }= (\bo{A} - \bo{L} \bo{C})\bo{e} 
\end{equation}

\end{flushleft}
\end{frame}



\begin{frame}{Observer gains, 1}
\begin{flushleft}

The observer $\dot {\bo{e} }= (\bo{A} - \bo{L} \bo{C})\bo{e} $ is \emph{stable} (i.e., the state estimation error tends to zero), as long as the following matrix has eigenvalues with negative real parts:

\begin{equation}
\bo{A} - 
\bo{L} \bo{C} \in \mathbb{H}
\end{equation}

We need to find $\bo{L}$. Let us observe a slight difference between observer design and controller design:

\bigskip

\begin{itemize}
    \item Controller design: find such $\bo{K}$ that $\bo{A} - \bo{B} \bo{K} \in \mathbb{H}$.
    \item Observer design: find such $\bo{L}$ that: $\bo{A} - \bo{L} \bo{C} \in \mathbb{H}$
\end{itemize}

\bigskip

The gain is on the left side for the observer, preventing the direct application of our design tools (pole placement, LQR) to tune it.

\end{flushleft}
\end{frame}


\begin{frame}{Observer gains, 2}
\begin{flushleft}

If $\bo{A} - \bo{L} \bo{C}\in \mathbb{H}$, then $(\bo{A} - 
\bo{L} \bo{C})^{\top}\in \mathbb{H}$ (eigenvalues of a matrix and its transpose are the same, see Appendix). 

\bigskip

Therefore, we can solve the following \emph{dual problem}:

\begin{itemize}
    \item find such $\bo{L}$ that $\bo{A}^{\top} - 
\bo{C}^{\top} \bo{L}^{\top} \in \mathbb{H}$.
\end{itemize}

\bigskip

This dual problem is \emph{equivalent} to a standard control design problem. We can solve it via pole placement or through the algebraic Riccati equation (LQR). In pseudo-code it can be represented the following way:
%
\begin{align*}
	\bo{L}^{\top} \texttt{= lqr}(\bo{A}^{\top}, \bo{C}^{\top}, \mathbf Q, \mathbf R)
	\ \ \ \text{or} \ \ \
	\bo{L}^{\top} \texttt{= place}(\bo{A}^{\top}, \bo{C}^{\top}, \mathbf p)
\end{align*}



where $\mathbf Q$ and $\mathbf R$ determine the speed and noise sensitivity of the observer.


\end{flushleft}
\end{frame}




\begin{frame}{Observer + Controller, 1}
\begin{flushleft}

Thus we get dynamics+observer combination:

\begin{equation}
\begin{cases}
\dot {\bo{x}} = \bo{A} \bo{x} + \bo{B} \bo{u} \\
\dot{\hat{\bo{x}}}  = \bo{A} \hat{\bo{x}} + \bo{B} \mathbf u + \bo{L}(\mathbf y - \bo{C} \hat{\bo{x}})\\
\bo{y} = \bo{C} \bo{x} \\
\bo{u} = -\bo{K} \hat{\bo{x}}
\end{cases}
\end{equation}

\bigskip

where $\bo{A} - \bo{B} \bo{K} \in \mathbb{H}$ and $\bo{A}^{\top} - 
\bo{C}^{\top} \bo{L}^{\top} \in \mathbb{H}$.

Let us re-write the dynamics:

\begin{equation}
	\begin{cases}
		\dot {\bo{x}} =\bo{A} \bo{x} - \bo{B} \bo{K} \hat{\bo{x}} 
		\\
		\dot{\hat{\bo{x}}}  = \bo{A} \hat{\bo{x}} - \bo{B} \bo{K}\hat{\bo{x}} +\bo{L}\bo{C} \bo{x} - \bo{L}\bo{C} \hat{\bo{x}}
	\end{cases}
\end{equation}


\end{flushleft}
\end{frame}




\begin{frame}{Observer + Controller, 2}
\framesubtitle{Stability analysis}
\begin{flushleft}

Let us re-write the dynamics

\begin{equation}
\begin{cases}
\dot {\bo{x}} = \textcolor{mydarkblue}{\bo{A}} \bo{x} \textcolor{mydarkpink}{- \bo{B} \bo{K}} \hat{\bo{x}} 
\\
\dot{\hat{\bo{x}}} = \textcolor{myyellow}{\bo{A}} \hat{\bo{x}} \textcolor{myyellow}{- \bo{B} \bo{K}} \hat{\bo{x}} + \textcolor{mydarkgreen}{\bo{L}\bo{C}} \bo{x} \textcolor{myyellow}{- \bo{L}\bo{C}} \hat{\bo{x}}
\end{cases}
\end{equation}

in a matrix form:

\begin{equation}
\begin{bmatrix}
\dot {\bo{x}} \\
\dot{\hat{\bo{x}}} 
\end{bmatrix}
=
\begin{bmatrix}
\textcolor{mydarkblue}{\bo{A}} & \textcolor{mydarkpink}{-\bo{B}\bo{K}}\\
\textcolor{mydarkgreen}{\bo{L}\bo{C}} & (\textcolor{myyellow}{\bo{A} - \bo{B}\bo{K}-\bo{L}\bo{C}})
\end{bmatrix}
\begin{bmatrix}
\bo{x} \\
\hat{\bo{x}}
\end{bmatrix}
\end{equation}

\bigskip

The eigenvalues of this matrix are not immediately obvious. But the problem can be simplified with a change of variables.

\end{flushleft}
\end{frame}




\begin{frame}{Closed-loop system}
\framesubtitle{Change of variables}
\begin{flushleft}

Let us use the following substitution: $\bo{e} = \bo{x} - \hat{\bo{x}}$, which implies $\hat{\bo{x}} = \bo{x} - \bo{e}$:

Our system had form:

\begin{equation}
\begin{cases}
\dot {\bo{x}} = \textcolor{mydarkblue}{\bo{A} \bo{x} - \bo{B}\bo{K} \hat{\bo{x}}} \\
\dot{\hat{\bo{x}}}  = \textcolor{mydarkpink}{\bo{A} \hat{\bo{x}} - \bo{B}\bo{K} \hat{\bo{x}} + \bo{L}\bo{C} \bo{x} - \bo{L}\bo{C} \hat{\bo{x}}}
\end{cases}
\end{equation}

Since $\dot{\bo{e}} = \textcolor{mydarkblue}{\dot{\bo{x}}} - \textcolor{mydarkpink}{\dot{\hat{\bo{x}}}}$, we get:
%
\[
\dot{\bo{e}} = 
\textcolor{mydarkblue}{\bo{A} \bo{x} - \bo{B}\bo{K} \hat{\bo{x}}} - 
\textcolor{mydarkpink}{(\bo{A} \hat{\bo{x}} - \bo{B}\bo{K} \hat{\bo{x}} + \bo{L}\bo{C} \bo{x} - \bo{L}\bo{C} \hat{\bo{x}})}
\]
%
\[
\dot{\bo{e}} = 
\bo{A} (\bo{x} - \hat{\bo{x}})  - \bo{L}\bo{C}(\bo{x} - \hat{\bo{x}})
\]
%
\[
\dot{\bo{e}} = 
(\bo{A}  - \bo{L}\bo{C})\bo{e}
\]

Equation $\dot {\bo{x}} = \bo{A} \bo{x} - \bo{B}\bo{K} \hat{\bo{x}}$ takes form:

\[
\dot {\bo{x}} = (\bo{A}-\bo{B}\bo{K}) \bo{x} +  \bo{B}\bo{K}\bo{e}
\]


\end{flushleft}
\end{frame}




\begin{frame}{Closed-loop system}
\framesubtitle{Upper triangular form}
\begin{flushleft}

Collecting $\dot {\bo{x}}$ and $\dot{\bo{e}}$ we get:

\begin{equation}
\begin{cases}
\dot {\bo{x}} = (\bo{A}-\bo{B}\bo{K}) \bo{x} +  \bo{B}\bo{K}\bo{e} \\
\dot{\bo{e}} = 
(\bo{A}  - \bo{L}\bo{C})\bo{e}
\end{cases}
\end{equation}

In matrix form it becomes:

\begin{equation}
\begin{bmatrix}
\dot {\bo{x}} \\
\dot{\bo{e}}
\end{bmatrix}
=
\begin{bmatrix}
(\bo{A}-\bo{B}\bo{K}) & \bo{B}\bo{K} \\
0 & (\bo{A}  - \bo{L}\bo{C})
\end{bmatrix}
\begin{bmatrix}
\bo{x} \\
\bo{e}
\end{bmatrix}
\end{equation}

Eigenvalues of an upper block-triangular matrix equal to the union of the eigenvalues of the blocks on the main diagonal (see Appendix B). Hence here, the eigenvalues of the system are equal to the union of eigenvalues of $(\bo{A}-\bo{B}\bo{K})$ and $(\bo{A}  - \bo{L}\bo{C})$. 

\end{flushleft}
\end{frame}



\begin{frame}{Separation principle}
\begin{flushleft}
 
Since the eigenvalues of the system are equal to the union of eigenvalues of $(\bo{A}-\bo{B}\bo{K})$ and $(\bo{A}  - \bo{L}\bo{C})$, we can make the following observation:

\bigskip

\begin{alertblock}{Separation principle}
As long as the observer and the controller are stable independently, the overall system is stable too.
\end{alertblock}

\bigskip

Therefore, the controller gain $\bo{K}$ and the observer gain $\bo{L}$  can be designed independently.

\end{flushleft}
\end{frame}






\begin{frame}{Further reading}
	
	\begin{itemize}
		
		
		\item Dorf \& Bishop, Modern Control Systems (Chapter 11.4 Observer Design; Chapter 11.5 Integrated Full-State Feedback and Observer)		
		
		
	\end{itemize}
	
\end{frame}

\myqrframe



\begin{frame}{Appendix A. Eigenvalues of transpose}
	\begin{flushleft}
		
		Given matrix $\bo{M}$ and its eigenvalue $\lambda$ and eigenvector $\bo{v}$. we can prove that $\lambda$ is an eigenvalue of $\bo{M}\T$.
		
		\begin{align}
			\bo{M}\bo{v} = \lambda \bo{v} \\
			\text{det}\ (\bo{M}- \bo{I} \lambda) = 0 \\
			\text{det}\ (\bo{M}\T- \bo{I} \lambda) = 0 \\
			\exists \bo{u} \ \text{s.t.} \ \bo{M}\T\bo{u} = \lambda \bo{u}
		\end{align}
		
		We used the fact that determinant of a matrix is equal to the determinant of its transpose: $\text{det}(\bo{A}) = \text{det}(\bo{A}\T)$.
		
	\end{flushleft}
\end{frame}


\begin{frame}{Appendix B, 1}
		\framesubtitle{Eigenvalues of block-triangular matrices}
	\begin{flushleft}
		
		Given matrix $\bo{M}$:
		%
		\begin{align}
			\bo{M} = 
			\begin{bmatrix}
				\bo{A} & \bo{B} \\ \bo{0} & \bo{C}
			\end{bmatrix}
		\end{align}
		
		Let $\lambda$, $\bo{v}$ be an eigenvalue and eigenvector of $\bo{A}$ and $\mu$, $\bo{u}$ be an eigenvalue and eigenvector of $\bo{C}$. We can prove that $\lambda$, $\bo{v}_M = \begin{bmatrix}
			\bo{v} \\ 0
		\end{bmatrix}$ are eigenvalue and eigenvector of $\bo{M}$:
	
		\begin{align}
			\begin{bmatrix}
				\bo{A} & \bo{B} \\ \bo{0} & \bo{C}
			\end{bmatrix}
		\begin{bmatrix}
			\bo{v} \\ \bo{0}
		\end{bmatrix}
	=
	\begin{bmatrix}
		\bo{A}\bo{v} \\ \bo{0}
	\end{bmatrix}
	=
	\lambda
	\begin{bmatrix}
		\bo{v} \\ \bo{0}
	\end{bmatrix}.
		\end{align}
		
	\end{flushleft}
\end{frame}



\begin{frame}{Appendix B, 2}
	\framesubtitle{Eigenvalues of block-triangular matrices}
	\begin{flushleft}
		
		If $\mu$ is not an eigenvalue of $\bo{A}$, we can prove that $\mu$, $\bo{u}_M = \textcolor{mydarkpink}{\begin{bmatrix}
				(\bo{I}\mu - \bo{A})^{-1} \bo{B} \bo{u}  \\  \bo{u} 
		\end{bmatrix}}$ are eigenvalue and eigenvector of $\bo{M}$:
	
		\begin{align*}
	\begin{bmatrix}
		\bo{A} & \bo{B} \\ \bo{0} & \bo{C}
	\end{bmatrix}
\begin{bmatrix}
	(\bo{I}\mu - \bo{A})^{-1}\bo{B} \bo{u}  \\  \bo{u} 
\end{bmatrix}
	=
	\begin{bmatrix}
		\bo{A}(\bo{I}\mu - \bo{A})^{-1}\bo{B} \bo{u} + \bo{B} \bo{u}
		 \\ 
		 \bo{C}\bo{u}
	\end{bmatrix}
	= \\
	=
	\begin{bmatrix}
		(\bo{I}+\bo{A}(\bo{I}\mu - \bo{A})^{-1}) \bo{B} \bo{u}
		\\ 
		\mu \bo{u}
	\end{bmatrix} 
=
	\begin{bmatrix}
	(\bo{I}\mu - \bo{A}+\bo{A})(\bo{I}\mu - \bo{A})^{-1} \bo{B} \bo{u}
	\\ 
	\mu \bo{u}
	\end{bmatrix} 
= \\
=
	\begin{bmatrix}
	\mu(\bo{I}\mu - \bo{A})^{-1} \bo{B} \bo{u}
	\\ 
	\mu \bo{u}
	\end{bmatrix} 
=
\mu 
\textcolor{mydarkpink}{
\begin{bmatrix}
(\bo{I}\mu - \bo{A})^{-1} \bo{B} \bo{u}
\\ 
\bo{u}
\end{bmatrix} 
}.
	\end{align*}	
	
	
	
	Counting the number of eigenvalues, we observe that eigenvalues of $\bo{M}$ include only eigenvalues of $\bo{A}$ and $\bo{C}$.
		
	\end{flushleft}
\end{frame}


\end{document}
